\section{System Architecture}
\label{sec:architecture}

GestureCAD adopts a decoupled client-edge architecture designed to balance real-time computer vision inference throughput with interactive 60~FPS graphics rendering. As illustrated in Figure~\ref{fig:architecture}, the system separates high-frequency optical tracking from the interactive geometric modeling client across a lightweight local WebSocket network layer.

\begin{figure*}[t]
    \centering
    \fbox{%
        \parbox[c][5.8cm][c]{0.92\linewidth}{%
            \centering\small\sffamily
            \textbf{Figure 2: System Architecture Diagram Placeholder}\\
            \vspace{0.3cm}
            \textit{Intended Diagram:}\\
            Detailed system topology showing dataflow across the two primary tiers:\\
            (1) \textbf{Edge Vision Service (Python)}: Commodity Webcam $\to$ OpenCV Capture Thread $\to$ MediaPipe Hands Inference (21 Landmarks) $\to$ Kinematic Signal Conditioning (Anatomical Normalization, Hysteresis FSM, Deadzone Gating, EMA) $\to$ JSON Telemetry Serializer $\to$ FastAPI WebSocket Server.\\
            (2) \textbf{WebSocket Transport}: Full-duplex JSON streaming on \texttt{ws://localhost:8000/ws/tracking}.\\
            (3) \textbf{Browser CAD Engine (React 18 + Three.js)}: Dual-path execution splitting into:\\
            -- \textit{High-Frequency Render Path}: \texttt{useFrame()} direct typed array (\texttt{Float32Array}) buffer mutations bypassing Virtual DOM $\to$ WebGL GPU Viewport;\\
            -- \textit{Transactional CAD State Path}: Ray-Plane Unprojection $\to$ Snapping \& Regularization $\to$ Zustand Store $\to$ Command History (Undo/Redo) $\to$ Geometry Engine (\texttt{ExtrusionEngine}, \texttt{RevolveEngine}, \texttt{SolidMesh}, \texttt{FaceEnumerator}) $\to$ Exporters (OBJ, DXF, STL, JSON).\\
            \vspace{0.3cm}
            \footnotesize Expected file: \texttt{paper/figures/fig2_architecture.png}
        }%
    }
    \caption{The decoupled client-edge architecture of GestureCAD. Real-time vision processing executes in a dedicated Python background process, while the interactive CAD application and WebGL rendering engine execute client-side in the web browser, communicating via local WebSockets.}
    \label{fig:architecture}
\end{figure*}

\subsection{Edge Vision Microservice}
The vision tracking subsystem is implemented as a multi-threaded Python 3.10+ service utilizing OpenCV for video frame acquisition and Google MediaPipe Hands for articulated skeletal inference. The subsystem executes an asynchronous event loop powered by FastAPI and Uvicorn:

\begin{enumerate}
    \item \textbf{Asynchronous Frame Acquisition}: A dedicated capture thread polls the host video device (defaulting to camera index 0) via \texttt{cv2.VideoCapture}. In empirical testing on host hardware, the acquisition loop achieved an average capture rate of $26.68\text{ FPS}$ with a standard $720\text{p}$ ($1280 \times 720$) resolution sensor.
    \item \textbf{Hand Pose Estimation}: Each captured frame is converted from BGR to RGB color space and processed by MediaPipe Hands configured with \texttt{static\_image\_mode=False}, \texttt{max\_num\_hands=1}, and \texttt{model\_complexity=1}. The neural network outputs 21 landmarks in normalized camera space $\mathbf{l}_i = (x_i, y_i, z_i) \in [0, 1]^3$, where $x$ and $y$ represent image plane coordinates and $z$ represents relative depth relative to the wrist joint.
    \item \textbf{Kinematic Feature Extraction}: The raw landmarks are ingested by the \texttt{GestureRecognizer} module, which evaluates anatomical scale baselines, joint angles, finger extension ratios, and the normalized pinch metric (formally derived in Section~\ref{sec:methodology}).
    \item \textbf{JSON Telemetry Serialization}: The resulting state is packaged into a compact JSON telemetry frame containing the discrete gesture state (\texttt{IDLE}, \texttt{OPEN\_PALM}, \texttt{FIST}, \texttt{PINCH}, \texttt{PEACE}), the toggle lock state, 2D filtered cursor coordinates, palm roll/swipe angles, the normalized pinch ratio, and high-resolution Unix timestamps. Telemetry frames are broadcast over a local WebSocket server at \texttt{ws://localhost:8000/ws/tracking}.
\end{enumerate}

\subsection{Browser CAD Client and Dual-Path Rendering}
The client-side modeling interface is implemented in React 18, utilizing Three.js and React Three Fiber (R3F) for WebGL scene graph management.

A critical design challenge in browser-based spatial interaction is handling streaming coordinate data without inducing frame drops. When high-frequency spatial streams ($30$--$60\text{ Hz}$) are routed through declarative React component state, the frequent state updates trigger continuous Virtual DOM reconciliation, leading to heap thrashing and garbage collection (GC) stutter.

To eliminate this bottleneck, GestureCAD implements a \textbf{dual-path execution architecture}:

\begin{enumerate}
    \item \textbf{High-Frequency Render Path (Bypassing Virtual DOM)}: Continuous tracking data (cursor coordinates, real-time stroke preview lines, and plane rotation handles) are maintained in mutable JavaScript references (\texttt{useRef}) rather than reactive component state. In the Three.js render loop (\texttt{useFrame}), the engine mutates the underlying WebGL \texttt{BufferGeometry.attributes.position.array} (\texttt{Float32Array}) directly in GPU-allocated memory and sets \texttt{needsUpdate = true}. This direct buffer mutation ensures continuous 60 FPS viewport feedback with zero Virtual DOM reconciliation overhead during active sketching.
    \item \textbf{Transactional CAD State Path (Command Pattern)}: Discrete geometric operations---such as committing a finished 2D profile, extruding a solid, executing a face selection, or performing Undo/Redo---are treated as discrete transactions. These operations are dispatched to a centralized Zustand state store and managed via a formal Command Pattern (\texttt{CommandManager.js}). Each command maintains inverse operations for deterministic Undo/Redo reversibility, updating the scene graph only upon transaction completion.
\end{enumerate}

\subsection{Coordinate Spaces and Transformation Pipeline}
GestureCAD coordinates geometric interactions across four distinct coordinate reference frames:
\begin{itemize}
    \item \textbf{Camera Normalized Space $\mathcal{S}_{\text{cam}}$}: Hand landmarks $\mathbf{l}_i \in [0, 1]^2$ normalized to webcam image dimensions.
    \item \textbf{Normalized Device Coordinates (NDC) $\mathcal{S}_{\text{ndc}}$}: Viewport screen coordinates mapped to $[-1, 1]^3$ for WebGL camera projection.
    \item \textbf{World Space $\mathcal{S}_{\text{world}}$}: Global three-dimensional Cartesian coordinate frame $\mathbb{R}^3$ wherein Three.js scene cameras, lighting, and instantiated CAD objects reside.
    \item \textbf{Local Working-Plane Space $\mathcal{S}_{\text{plane}}$}: 2D parametric coordinate space $(u, v) \in \mathbb{R}^2$ established on an active working plane (XY, YZ, XZ) or dynamically attached to an existing polyhedron face via an orthonormal transformation $\mathbf{R} \in SO(3)$.
\end{itemize}
The mathematical mappings connecting these coordinate spaces are detailed in the subsequent section.
